\documentclass[10pt,a4paper]{amsart}
\usepackage[margin=19mm]{geometry}
\usepackage{amsmath,amssymb,mathtools}
\usepackage[scr=rsfs,cal=euler]{mathalfa}
\usepackage{xfrac}
\usepackage[unicode,hidelinks,bookmarks=false]{hyperref}
\newcommand{\ie}{i.e.\ }
\newcommand{\cf}{cf.\ }
\newcommand{\resp}{respectively\ }
\newcommand{\Subsection}{\subsection}
\providecommand{\nobreakdash}{\nobreak\hbox{-}}
\setlength{\emergencystretch}{2em}
\multlinegap0pt
\usepackage{enumerate}
\usepackage{amssymb}
\usepackage{float}
\usepackage{mathtools}
\usepackage{xcolor}
\newtheorem{lemma}{Lemma}
\title{\textbf{On the Equivalence of Zero-Sum Games and Conic Programs}}
\author{Nikos Dimou}
\date{Source: arXiv:2302.03066v5 (CC BY 4.0)}
\begin{document}
\maketitle

\noindent\emph{Source and licence.} \textbf{On the Equivalence of Zero-Sum Games and Conic Programs}. Version arXiv:2302.03066v5 (CC BY 4.0), distributed by arXiv under the Creative Commons Attribution 4.0 International licence, http://creativecommons.org/licenses/by/4.0/, as recorded in the arXiv licence metadata for this exact version and shown on the abstract page https://arxiv.org/abs/2302.03066v5. Full licence notice: \texttt{licenses/arxiv-2302.03066-CC-BY-4.0.txt}.\par\smallskip

\noindent\textit{Editorial scope.} Counted unit: the article's lemma lem:3.6 (source0.tex lines 567--569). The article's complete original proof of it is delivered with it (source0.tex lines 570--591, one \texttt{proof} environment of 22 lines). No mathematics is written, repaired or re-derived: the statement and the proof are the article's own lines, in the article's own notation, and no retained line is substituted or re-flowed.  Retained uncounted as the source-local context the proof needs: lines 324--326.  Nothing was omitted from any retained span, so the package carries no omission record.  No retained line contains a citation, so the wrapper carries no bibliography and no bibliography entry.  No retained line contains a figure environment, an image-inclusion command or drawing code, so no image is delivered and the package declares no asset dependency.  Editorial formatting only, in addition: an AMS \texttt{amsart} preamble that restates the article's own declarations and macros used by the retained lines, namely the environments \texttt{lemma} on separate counters, together with the standard packages those lines need.  The retained environments print in this document as Lemma 1 in order of appearance, whereas the article prints the same items with the numbering of its own full document.  The complete native source of the article as distributed in the arXiv e-print for this exact version is bundled unchanged as \texttt{sources/137052/source0.tex}.
\begin{equation}\label{eq:2.6}
    u(x^*,y)\geq u(x^*,y^*)\geq u(x,y^*)\;\;\forall x\in S,\;\forall y\in T.
\end{equation}

\begin{lemma}\label{lem:3.6}
     Let $\alpha\in \text{int}C^*$, $\beta\in \text{int}K^*$, and a linear operator $A:X\rightarrow Y^*$. Consider the linear operators $E:X\rightarrow Y^*$ with $E(x):= \beta\langle\alpha,x\rangle$, and $B:X\rightarrow Y^*$ with $B(x):=\lambda A(x)+\kappa E(x)$ for some fixed $\lambda,\kappa >0$. Also, let $v_A$ be the game value for $G_A=(\alpha,\beta,A)$ and $v_B$ be the corresponding one for $G_B=(\alpha,\beta,B)$, where both games are defined by \hyperlink{f1}{(F1)}-\hyperlink{f3}{(F3)}. Then, $v_A$ exists if and only if $v_B$ exists, and in that case $v_B=\lambda v_A+\kappa$. Moreover, $(x^*,y^*)$ is a Nash equilibrium for $G_A$ if and only if it is a Nash equilibrium for $G_B$.
\end{lemma}

\begin{proof}
As the strategy sets are bases of convex cones, the following relation holds:
\begin{equation}\label{eq:B}
    \langle y,Bx\rangle=\lambda\langle y,Ax\rangle +\kappa\langle y,\beta\rangle\langle \alpha,x\rangle=\lambda\langle y,Ax\rangle+\kappa\;\;\forall x\in S,\;\;\forall y\in T.
\end{equation}
Let $\underline{v}_A$ be the lower value for $G_A$ and $\underline{v}_B$ be the corresponding one for $G_B$. Then, (\ref{eq:B}) yields
\begin{equation*}
    \underline{v}_B=\sup_{x\in S}\inf_{y\in T}\langle y,Bx\rangle=\lambda\sup_{x\in S}\inf_{y\in T}\langle y,Ax\rangle +\kappa =\lambda\underline{v}_A+\kappa.
\end{equation*}
Similarly, if $\overline{v}_A,\;\overline{v}_B$ are the upper values for $G_A$ and $G_B$, respectively, then $\overline{v}_B=\lambda\overline{v}_A+\kappa$. Therefore, $v_B$ exists if and only if $v_A$ exists, and the relation $v_B=\lambda v_A+\kappa$ follows.\par
For the second part, let $(x^*,y^*)$ be a Nash equilibrium for $G_A$. Then, by (\ref{eq:2.6}) and (\ref{eq:B}) we equivalently have
\begin{equation*}
    \langle y,Ax^*\rangle\geq\langle y^*,Ax^*\rangle\geq\langle y^*,Ax\rangle\;\;\forall x\in S,\;\;\forall y\in T\;\Leftrightarrow
\end{equation*}
\begin{equation*}
    \lambda\langle y,Ax^*\rangle+\kappa\geq\lambda\langle y^*,Ax^*\rangle+\kappa\geq\lambda\langle y^*,Ax\rangle+\kappa\;\;\forall x\in S,\;\;\forall y\in T\;\Leftrightarrow
\end{equation*}
\begin{equation*}
    \langle y,Bx^*\rangle\geq\langle y^*,Bx^*\rangle\geq\langle y^*,Bx\rangle\;\;\forall x\in S,\;\;\forall y\in T,
\end{equation*}
which is equivalent to $(x^*,y^*)$ being a Nash equilibrium for $G_B$.
\end{proof}

\end{document}
